\chapter{Conclusions}
\label{cap:conclusions}

In conclusion, GeoMedia presents an open-source, privacy-oriented and extensible platform for location-based sharing of multimedia content, including audio, text, images and other file types.

Designed as a flexible framework, the application enables users to build their own collections and define different ways to share content and interact with posts present in the environment surrounding them. This approach supports a variety of use cases while providing users with a more contextual and personalized interaction with geolocated media. 
Furthermore, GeoMedia integrates a range of distinctive functionalities that differentiate it from conventional social networking platforms by adopting a content-centric rather than profile-centric approach, in which posts and their spatial context constitute the core elements of the social experience. This design also aims to reduce some of the social pressures commonly associated with profile-centric platforms, such as social comparison and FoMO, by shifting the focus from users and their social status towards shared content and experiences.
Overall, GeoMedia demonstrates the potential of geolocated content as a means for digital connection with the physical world and social engagements. Its open, extensible and modular architecture, providing the foundations for further development and experimentation in areas such as storytelling, interactive location-based games and place discovery. GeoMedia can be considered not only a platform for sharing media, but also as a geolocated media hub that explores how technology can foster meaningful interactions between people, digital content and their surrounding environment.

